\section{Discussion}
\label{sec:discussion}

The selection comparison uses three inputs: an entrance phase with
a bounded correction and a uniform vanishing remainder, a same-index
return error smaller
than $1/\log(1/\eps)$, and localization of the actual stopping
event to a narrower window. Together they explain why parameter
dependence disappears for a free target but determines on which
side of a contact the first crossing occurs. The abstract criterion
in Remark~\ref{sel:criterion} records this last comparison;
the original preparation, admissible return domains and sharp
full-state budget establish its hypotheses here. For complete
paths one also needs convergence of the incoming motion, ordered
comparison of any fixed number of final returns, and a common
shrinking neighborhood for the whole earlier tail.

The normal filter, the $3/4$ return exponent, the $9/32$ entrance
exponent and the signs ordering $\chi$ come from the specified
mechanical and constitutive laws. Compactly supported profile
variations share an exact incoming tail, so their critical-passage
response reduces to a finite-time variational problem
(Lemma~\ref{lem:profile-variation}). Controllability and the
single-visit high-slip interval then supply a direction with nonzero
response, with a common mechanical box and section chosen for the
constructed compact curve. The profile tails and mechanical rate
margins also enter the preparation estimates. The prescribed compact $\beta$ range makes the preparation
estimates uniform; it is the range treated by the present proof.
The signed first contact permits $r=O(\delta)$, whereas the
preparatory flights put later inputs in $r\le A\delta^{4/3}$.
A different detuning scale or normal law would require its own
entrance and return estimates before the selection comparison
could be applied.

The observation results distinguish a longer path from the first
fixed-angle prefix that it contains. A strict local crossing gives
continuous prefix selection, while the angle section contains the
whole limiting contact and has a discontinuous first-hit operation.
The additional scaled-memory and path-identification results retain
information that can be lost by projecting the state or discarding
the contact order.

The selection results compare fixed labels. Whether belt or profile
parameter differences of order $\eps\log(1/\eps)$ yield a transition
law remains open. The absolute shadowing estimate used here does
not control such parameter increments.
